%% P0: The Invisible Civilization
%% Draft v0.4 — Tier 2 drafting: §4-§9 added (Selective Survival, Wayang, Six-Layer Framework,
%%             Predictions, Limitations, Conclusions). Builds on v0.3 (decisive critiques closed
%%             per DeepSeek re-review 2026-04-22).
%% Target venue: Journal of Anthropological Archaeology (Elsevier, Q1, subscription route)
%% Target length: 10,000-12,000 words total; this file targets ~10,500 words
%% Companion to: P1-core (sedimentation calibration, submitted separately to JASREP)
%% Lineage: v0.1 (2026-04-20, §1-3.2) → v0.2 (2026-04-22, +§3.3-§3.6) →
%%          v0.3 (2026-04-22, decisive-fixes §3.4 + §3.6 + §3.2 paradox + §2.2 sensitivity) →
%%          v0.4 (2026-04-22, +§4-§9 full paper)

\documentclass[12pt,a4paper]{article}
\usepackage[utf8]{inputenc}
\usepackage[T1]{fontenc}
\usepackage{graphicx}
\usepackage{amsmath}
\usepackage{natbib}
\usepackage{booktabs}
\usepackage{array}
\usepackage{hyperref}
\usepackage[left=2.5cm,right=2.5cm,top=2.5cm,bottom=2.5cm]{geometry}
\usepackage{lineno}
\usepackage{setspace}
\doublespacing
\linenumbers

\hypersetup{colorlinks=true,linkcolor=blue,citecolor=blue,urlcolor=blue}

\begin{document}

\title{The Invisible Civilisation:\\Six Independent Lines of Evidence for an Archaeologically Erased Pre-Hindu Nusantara}

\author{Mukhlis Amien\textsuperscript{1,*} \and Go Frendi Gunawan\textsuperscript{1}}

\date{}

\maketitle

\noindent\textsuperscript{1}Lab Data Sains, Universitas Bhinneka Nusantara, Malang, Indonesia\\
\textsuperscript{*}Correspondence: amien@ubhinus.ac.id

\begin{abstract}
The archaeological record of Java before the fourth century~CE is effectively empty. The conventional narrative treats this as a chronological fact: Javanese polities formed late and adopted Indic inscription with Indic civilisation. Demographic modelling, ecological carrying capacity estimates, and comparison with contemporaneous Austronesian societies suggest a different possibility: Java at 400~CE supported roughly one to two million inhabitants, comparable to Dvaravati Thailand and exceeding the early Philippines and Sriwijaya. If this population existed, its material record has been selectively erased. This paper assembles six independent evidence channels that converge on the existence of such an erased substrate civilisation: cumulative volcanic sedimentation rates that place pre-Hindu deposits at 4--10~m depth; deep comparison with the Philippine archaeological record, which preserves 275--340 pre-400~CE sites (55--65\% open-air) against zero in Java's volcanic interior; linguistic substrate reconstruction identifying Proto-Austronesian writing vocabulary and hundreds of substrate items in six Sulawesi languages; reinterpretation of published paleogenomic data showing that ancient DNA has never been recovered from any open-air context in volcanic Java, a pattern consistent with the karst-versus-open-air preservation differential that also characterises Southeast Asian paleogenomic recovery more broadly; colonial-archive text mining that surfaces 22{,}000+ structured mentions matching predicted burial-depth gradients; and archaeometric evidence in two chronological layers---locally-produced Pejeng-type bronze drums dated to the 1st--3rd centuries~CE (pre-inscription anchor), and Javanese \textit{Jatim} glass beads in 4th--6th-century~CE elite tombs across China, Korea, and the Roman Red Sea (showing the invisibility regime extends into the early state period)---that together attest to an organised production system whose workshops remain archaeologically absent from Java and Bali while their outputs survive 4{,}500--8{,}000~km away. The six channels use different data, different methods, and different observer communities. They converge on a ``selective survival'' regime in which metal and cave-preserved material survives while organic settlements do not. We argue this regime is diagnosable computationally, falsifiable through specified subsurface predictions, and visible in living tradition (wayang, \textit{gunungan}, the Semar complex). The paper closes with a six-filter framework and eight pre-registered predictions whose outcomes can support, constrain, or falsify the argument.

\medskip
\noindent\textbf{Keywords:} archaeological taphonomy; multi-channel convergence; pre-Hindu Nusantara; invisible civilisations; selective survival; volcanic tropical archaeology; Indonesian prehistory
\end{abstract}


%% ================================================================
\section{Introduction: The Dwarapala's Half-Buried Body}

In 1803, the Dutch colonial administrator Nicolaus Engelhard encountered a striking scene at the ruins of Candi Singosari in East Java. Two enormous stone guardian statues, each 370~cm tall and carved from monolithic andesite in 1268~CE during the reign of King Kertanegara, stood with approximately half their bodies below the surrounding ground surface. They had not been deliberately interred. Gunung Kelud had erupted about twenty times in the intervening 535 years, and the landscape had risen around them at approximately 3.5~mm per year \citep{kinney2003}. The statues were visible only because of their monumental scale. Structures built of wood, brick, and thatch in the same period left no comparable marker above the accumulating sediment.

This observation opens a question that has not been pursued systematically in Indonesian archaeology. If a landscape buries stone monuments to half their height within five centuries, what does it do to the organic architecture that housed most of its inhabitants? The rate Engelhard's observation implies, when extended with three additional calibration sites from the Merapi volcanic system (Sambisari, Kedulan, Kimpulan; 9th century~CE), yields a cross-system mean of approximately 4.4~mm/yr \citep[hereafter P1-core]{amien_p1core}. Applied to the pre-Hindu horizon, this rate places archaeological deposits from approximately 400~CE at depths of 3.9--10.1~m in Java's volcanic basins---beyond the reach of surface survey and at or exceeding the effective limit of conventional ground-penetrating radar.

The calibration itself is geomorphological, not archaeological. It cannot demonstrate that pre-Hindu settlements exist at those depths. It can only establish that the surface record of volcanic Java has a detection horizon that cuts off sharply within the historical period. What lies below that horizon is the subject of this paper.

The conventional reading of Java's early archaeological record is developmental: Hindu-Buddhist kingdoms crystallised around the 4th century~CE, adopted Indic inscription technology, and left the epigraphic and monumental record from which subsequent polities are reconstructed \citep{coedes1968,miksic2004}. Before this threshold: silence. The silence has typically been interpreted as absence. Kutai Martadipura in non-volcanic East Kalimantan, dated to approximately 400~CE through its Yupa inscriptions \citep{vogel1918}, is treated as Indonesia's oldest known polity. Javanese state formation is dated from the 4th or 5th century~CE inscriptions of Tarumanagara. The possibility that a comparable or larger population existed on volcanic Java before these dates, but left no surface-recoverable material record, has been considered implicitly impossible---or perhaps, more often, simply not considered.

This paper argues that the possibility is both plausible and evidenced. Not through any single discovery, but through the convergence of six independent lines of evidence: sedimentation-rate calibration, comparative archaeology against the Philippine record, linguistic substrate reconstruction, reinterpretation of published paleogenomic data, colonial-archive text mining, and archaeometric evidence from exported Javanese material culture in extra-Javanese tomb and port contexts. Each channel addresses a different aspect of the problem; each uses data collected by a different research community; each was not designed with VOLCARCH's thesis in mind. What they share is a single pattern: volcanic Java's pre-400~CE material record is quantitatively incompatible with what Austronesian prehistory, ecological carrying capacity, linguistic depth, genetic continuity, and documented material exports would independently predict.

We develop the argument in three parts. Section~2 establishes the demographic puzzle: how many people should have lived on Java at 400~CE, and what archaeological footprint should they have left? Section~3 presents the six evidence channels, each with its own method, result, independence assessment, and limitation. Section~\ref{sec:selective_survival} introduces the concept of ``selective survival'' that organises the positive evidence---bronze drums, megaliths, wayang---alongside the negative. Section~\ref{sec:layers} develops a six-filter framework linking the evidence channels into a single cascade of archaeological invisibility. Section~\ref{sec:predictions} presents eight pre-registered predictions whose outcomes can falsify, support, or constrain the framework. Section~\ref{sec:limits} catalogues limitations honestly. Section~\ref{sec:conclusion} concludes with methodological implications for Island Southeast Asian archaeology.

A note on what this paper does not claim. We do not claim to have discovered a lost civilisation. We do not claim that any specific location contains pre-Hindu remains. We do not claim that the demographic estimates are precise---they span an order of magnitude, and we present them as bounds rather than measurements. We do claim that the five channels, when read together, make a purely developmental explanation of Java's pre-400~CE record quantitatively untenable, and that the alternative explanation---archaeological invisibility produced by compound taphonomic processes---is both testable and consistent with what positive evidence does survive.


%% ================================================================
\section{The Demographic Puzzle}
\label{sec:puzzle}

\subsection{What the record shows}

Java's pre-400~CE archaeological record is minimal. No inscriptions predate the 4th-century Tarumanagara stones in non-volcanic West Java \citep{vogel1918,coedes1968}. No monumental architecture predates the 7th-century Dieng temples in the Kedu plain. Open-air habitation sites identified as pre-Hindu exist in the low single digits, most with contested or imprecise dating, and all outside Java's central volcanic interior. A compilation of 666 known archaeological sites in East Java \citep[P1-core]{amien_p1core} yields fewer than five percent attributed to any period earlier than the 10th century~CE, and zero from Java's volcanic interior that can be confidently dated before 400~CE.

This record is concentrated in three ways. First, geographically: known sites cluster within 50~km of major volcanic centres in the Singosari-Majapahit heartland around Malang, Mojokerto, and the Prambanan plain. Second, materially: the record is dominated by stone architecture (\textit{candi}) that weighs thousands of tonnes and protrudes through accumulated sediment. Third, chronologically: it favours the 8th through 15th centuries~CE, the periods when Hindu-Buddhist polities commissioned stone construction. Each concentration reflects a different bias. Together they produce a record that maps monument durability, survey effort, and genre-specific inscriptional practice---not settlement distribution.

\subsection{What Austronesian ecology predicts}

A population estimate for 400~CE Java can be derived from three independent approaches, all of which converge on the same order of magnitude.

The ecological carrying capacity approach applies ethnographically documented densities for tropical Austronesian agricultural systems---approximately 5 to 30 persons per km\textsuperscript{2} for swidden and early wet-rice cultivation \citep{bellwood2017,kirch2000,baylisssmith1980}---to Java's habitable lowlands. Using a terrain-suitability threshold that excludes slopes steeper than roughly 15 degrees and elevations above 1{,}000~m, habitable area is approximately 65{,}000~km\textsuperscript{2}, yielding a carrying-capacity range of 325{,}000 to 1{,}950{,}000 inhabitants.

Demographic back-projection from early-modern anchors provides an independent check. Reid's compilation places Java's 1600~CE population at approximately four to five million \citep{reid1988}. Plausible Malthusian growth rates (0.05--0.15\% per annum averaged across the 1200 intervening years, accounting for eruption disruptions, trade-cycle fluctuations, and the 14th-century Black Death analogues) imply a 400~CE baseline on the order of 600{,}000 to two million---consistent with the carrying-capacity result without sharing any of its assumptions. We acknowledge the circularity risk: growth rates are not directly observable over the pre-modern period, and one reviewer of an earlier version of this paper correctly noted that faster growth driven by post-400~CE technological and institutional change (iron agriculture, wet-rice intensification, political centralisation) would imply a lower 400~CE baseline. Under a doubled growth rate (0.3\% per annum, a plausible upper bound for an Indianisation-accelerated trajectory) the back-projection falls to 150{,}000--250{,}000 inhabitants at 400~CE. This lower bound is still quantitatively incompatible with zero observed settlements in the volcanic interior: at 100--200 persons per village, 150{,}000 inhabitants imply between 750 and 1{,}500 settlements, a gap of two to three orders of magnitude against the observed count.

Comparative scaling against contemporaneous Austronesian societies places Java at or above the regional median. Dvaravati Thailand supported 300{,}000--500{,}000 inhabitants during roughly the same period \citep{higham2014}; the early Philippines 200{,}000--500{,}000 \citep{junker1999}; early Sriwijaya 50{,}000--150{,}000 \citep{manguin2004}. Java, with the densest volcanic-soil fertility in the archipelago, would be an outlier if its 400~CE population fell below the Thai figure.

A fourth, non-demographic data point constrains the lower bound. The \textit{Hou Han Shu} records an embassy to Han China in 132~CE from the king of Ye-tiao, identified on Middle Chinese phonological grounds as Yavadvipa, that is, Java \citep{pelliot1916,wolters1967}. Whatever its population, the polity behind that embassy was organised enough to dispatch a diplomatic mission to an empire 4{,}500~km away. This predates the conventional Kutai Yupa dating by roughly two and a half centuries and implies that whatever demographic and institutional base sustained an international embassy already existed in Java in the 2nd century~CE. It does not settle the population figure, but it raises the minimum organisational threshold the figure must accommodate.

The three approaches converge on a central estimate of one to two million inhabitants at 400~CE, with a conservative lower bound near 600{,}000 and an upper bound near three million. This convergence is not coincidental; the approaches share no data and do not use each other's intermediate quantities. We adopt 1--2 million as a working figure, with sensitivity to the lower bound tested below.

\subsection{What the gap implies}

At standard Austronesian village sizes of 100 to 200 persons, a population of one to two million implies between 5{,}000 and 20{,}000 settlements. The observed count of pre-400~CE settlement sites in volcanic Java's interior is zero to three, depending on dating confidence. The discrepancy, computed as expected settlements divided by observed sites, spans approximately three orders of magnitude under plausible parameter choices. Under the minimal-population scenario (590{,}000 inhabitants, maximum village size of 200, and liberal counting of 3 sites) the gap is approximately 1{,}000-fold. Under the moderate scenario (1{.}9 million inhabitants, same village and site parameters) the gap is approximately 3{,}200-fold (E108). Under the maximum scenario (3{.}9 million inhabitants, same parameters) the gap reaches approximately 6{,}500-fold. Across the full range of scenario parameters tested, the gap remains above 1{,}000-fold.

Three potential explanations present themselves. First, the population estimate could be wrong. Second, the settlements could have taken forms that leave no archaeological trace regardless of taphonomy. Third, the archaeological record could be biased against detecting them.

The first explanation is difficult to sustain. To close the gap through lower population alone would require 400~CE Java to have fewer than 2{,}000 inhabitants---a figure inconsistent with every independent estimation method, every comparative analogue, and every demographic back-projection that reaches the period. Java would need to be empty, in an archipelago where neighbouring Austronesian societies with far poorer soils supported hundreds of thousands.

The second explanation---that settlements left no archaeological trace---has surface plausibility. Organic construction using bamboo, wood, palm leaf, and thatch is undetectable after a few centuries in the tropical climate even without volcanic burial \citep{erasmus2005}. However, the organic-architecture explanation alone predicts that non-volcanic tropical Austronesian settings should also show similar archaeological absence. They do not. The Philippine record, discussed in Section~\ref{sec:channels}, preserves hundreds of open-air pre-400~CE sites across varied environments. The absence, in other words, is specifically concentrated in Java's volcanic terrain, not distributed across all comparable tropical Austronesian landscapes.

The third explanation---taphonomic bias in the archaeological record---is the hypothesis this paper develops. Not as a replacement for the first two (population was certainly smaller than the 20{,}000-settlement upper bound, and organic architecture certainly contributes to invisibility) but as the load-bearing factor that, together with the others, makes the discrepancy coherent. The remainder of this paper assembles the five evidence channels that support this interpretation.


%% ================================================================
\section{Six Independent Evidence Channels}
\label{sec:channels}

The argument proceeds through six channels, each drawing on different data, different methods, and different observer communities. We assess each separately, then discuss their combined implications. The channels are: cumulative volcanic sedimentation (\S\ref{sec:ch1}); comparative archaeology against the Philippine record (\S\ref{sec:ch2}); linguistic substrate reconstruction (\S\ref{sec:ch3}); paleogenomic reinterpretation (\S\ref{sec:ch4}); colonial archive text mining (\S\ref{sec:ch5}); and archaeometric evidence from exported Javanese material culture (\S\ref{sec:ch6}).

Table~\ref{tab:channels} summarises the channels, their primary evidence, and their independence from one another.

\begin{table}[t]
\caption{The six evidence channels, primary sources, and independence assessment. ``Independent of'' indicates channels that share no primary data or method. Note: all channels that invoke population or comparative estimates draw ultimately on the same ethnographic and historical-demographic literature base (principally \citealp{bellwood2017,reid1988,kirch2000,higham2014}); the channels are methodologically and evidentially diverse rather than fully epistemically independent. Channel 6, which relies on data collected by Chinese, Korean, and Roman-world specialists studying finds in their own regions, is the most fully independent of VOLCARCH's evidential base.}
\label{tab:channels}
\begin{tabular}{p{2.5cm}p{4.5cm}p{4cm}p{3cm}}
\toprule
Channel & Primary data & Observer community & Independent of \\
\midrule
Sedimentation (\S\ref{sec:ch1}) & 4 calibration sites + 51 eruption--site pairs & Geoarchaeologists, volcanologists & 2, 3, 4, 5, 6 \\
Philippines comparison (\S\ref{sec:ch2}) & ${\sim}275{-}340$ Philippine pre-400 CE sites & Southeast Asian archaeologists & 1, 3, 4, 6 (partial 5) \\
Linguistic (\S\ref{sec:ch3}) & ABVD wordlists, OJW Wordnet, 438 substrate items, DHARMA inscriptions & Historical linguists, computational linguists & 1, 2, 4, 6 (partial 5) \\
Genomic (\S\ref{sec:ch4}) & Published aDNA + WGS datasets & Paleogeneticists & 1, 2, 3, 5, 6 \\
Colonial archive (\S\ref{sec:ch5}) & 16 OV volumes, 1{,}768 Delpher articles & Colonial-era Dutch and Indonesian fieldworkers & 3, 4, 6 (partial 1, 2) \\
Archaeometric (\S\ref{sec:ch6}) & Pejeng bronze drums (1st--3rd c.~CE, pre-inscription) + \textit{Jatim} glass beads (4th--6th c.~CE, Datong, Gyeongju, Berenike) & Metallurgists, Chinese archaeometrists, Korean archaeologists, Roman Red Sea excavators & 1, 2, 3, 4, 5 \\
\bottomrule
\end{tabular}
\end{table}

\subsection{Volcanic sedimentation rates and the detection horizon}
\label{sec:ch1}

The first channel is the one whose methodological detail is developed in full in a companion calibration paper \citep{amien_p1core}. Here we summarise its result for the synthesis.

Four archaeological sites with independently documented construction dates and measured burial depths were compiled from two volcanic systems. The Dwarapala statues of Singosari (1268~CE, Kelud system, 185~cm of accumulation over 535 years) yield a rate of 3.5~mm/yr, internally consistent with Kelud's approximately twenty eruptions in that period. Three Merapi-system temples---Sambisari (835~CE, 500--650~cm), Kedulan (869~CE, 600--700~cm), Kimpulan (900~CE, 270--500~cm)---yield rates of 4.4--5.7, 5.3--6.2, and 2.4--4.5~mm/yr respectively. The four-site mean is approximately 4.4~mm/yr across 350~km of Java and two geologically distinct volcanic basins.

Two independent validation datasets test this rate. First, 51 eruption--archaeological site pairs, compiled from colonial and volcanological literature spanning 12 volcanic systems, yield a mean measured burial depth of 3.41~m. Applying the four-site calibration rate to a 760-year horizon predicts 3.3~m---an agreement within 4\% that was not a design target of either compilation. Second, a larger compilation of 363 site-depth observations from East Java archaeological service reports (P1-core, E075) yields a Pearson correlation of $r=0.951$ between predicted and observed burial depths.

Applied as a first-order linear projection, the calibration rates place pre-400~CE deposits in the Malang basin at 3.9--10.1~m depth, Kanjuruhan-era (760~CE) deposits at 3.0--7.9~m, and Singosari-era (1268~CE) deposits at 1.8--4.7~m. The effective range of ground-penetrating radar in volcanic sediment is 2--5~m \citep{conyers2004,goodman2013}; the effective range of systematic surface survey is below 1~m. The pre-400~CE horizon is therefore beneath both techniques' practical reach, except through construction-triggered exposure or deep coring.

The channel establishes a detection horizon, not a buried record. It cannot show that settlements exist at the projected depths; it can show that the surface archaeological record of Java's volcanic interior cuts off within the historical period. Whether the cut-off reflects taphonomic invisibility or cultural absence is a question the other four channels address.

This channel is falsifiable. Systematic deep coring or construction-triggered exposure to depths of 4--10~m in volcanic-interior target zones that yields zero anthropogenic material, across a statistically adequate sample, would reject the detection-horizon interpretation. Rates systematically lower than 4.4~mm/yr in additional calibration sites would reduce the projected pre-400~CE depth and invalidate the claim that this horizon exceeds conventional survey reach.

\subsection{The Philippines natural experiment}
\label{sec:ch2}

The Philippines shares Java's major cultural and environmental parameters: Austronesian heritage, tropical climate, island-archipelagic geography, active volcanism, and organic building tradition. It differs in three relevant ways: lower volcanic density (approximately 0.07 active volcanoes per 1{,}000~km\textsuperscript{2} versus Java's 0.35), abundant karst providing natural preservation contexts, and smaller pre-colonial population by all estimates. If volcanic burial were not a significant archaeological factor, the Philippine pre-400~CE record should be proportionally thinner than Java's given its smaller population; if it is, the record should be proportionally richer given its preservation advantages.

A deep composition analysis of the Philippine archaeological record \citep[E201]{amien_e201} identifies approximately 275--340 documented pre-400~CE sites across site types: around 30--50 cave habitation sites (Tabon, Callao, Balobok, Pilanduk, Bilat), approximately 170 open-air Paleolithic sites (Cagayan Valley alone documents 168+), 30+ Neolithic shell middens (the Lal-lo/Gattaran complex), 20--30 open Neolithic settlement sites (Andarayan, Batanes jade sites), 15--20 Metal Age cave burials, 5--10 Metal Age open burials, and 5--10 rock art sites. The site-type breakdown is 55--65\% open-air, 35--45\% cave-based---a diverse record, not a cave-dominated one.

Material culture representation follows the same pattern. The Philippine pre-400~CE record includes abundant pottery (from approximately 4{,}000~BP onward), stone tool assemblages (700{,}000~BP through Neolithic), metal objects (500~BCE onward, including gold, bronze, and iron), tens of thousands of jade artifacts from the Batangas culture (2000--1500~BCE), hundreds of human burials in both cave and open contexts, evidence of rice from 1500--1400~BCE (Andarayan), and domesticated pig remains from 2500--2200~BCE (Nagsabaran). Volcanic Java's pre-400~CE record contains none of these categories.

Critically, Philippine volcanic zones themselves preserve archaeological material. Batangas---adjacent to Taal volcano, which has erupted 39 or more times historically---preserves the Calatagan burial grounds, jade culture sites, and coastal shell middens. Batanes sites buried under Iraya ash were recovered by systematic excavation \citep{bellwood_dizon2013}. The site-forming potential of volcanic burial, demonstrated in the Philippines, is recovery-dependent rather than preservation-destroying.

Volcanic Java's zero pre-400~CE open-air sites stand against documented pre-400~CE sites in Philippine volcanic zones. Within Java itself, non-volcanic West Java preserves pre-Hindu material: the Buni culture (c.~400~BCE to 500~CE) on the northern coast and Batujaya's 1st--3rd~century~CE prehistoric burial layer beneath the later Tarumanagara temple complex demonstrate that open-air pre-400~CE archaeology is recoverable where volcanic sedimentation is absent \citep{manguin2011,wibisono1994}. The within-island contrast---non-volcanic Java preserves, volcanic Java does not---constrains the possible explanations. The remaining factors are (a) Philippines' more abundant karst, which provides natural preservation capsules bypassing other taphonomic filters; (b) Java's more continuous volcanic chain, producing more extensive tephra blankets; and (c) the assumption in Indonesian archaeology that the volcanic interior is empty, which has directed survey effort elsewhere and become self-fulfilling.

The Philippines comparison provides the strongest single external constraint on the Java thesis. It controls for climate, Austronesian heritage, tropical ecology, organic architecture tradition, and tropical organic decay. What remains as differential is volcanism, karst availability, and survey attention. The zero-versus-hundreds site-count contrast is not reducible to any single one of these factors alone.

A question naturally follows: if a substantial pre-400~CE interior population existed in volcanic Java, why is its trade-goods footprint not more visible in the neighbouring non-volcanic zones (Buni, Batujaya) where preservation is demonstrably better? The question deserves an explicit answer. Three points bear on it. First, the coastal West Java sites \textit{do} contain material-culture assemblages (pottery, iron, glass beads, bronze) indicating participation in wider Island Southeast Asian exchange networks that necessarily included some interior-to-coastal transfer---the absence at issue is specifically the absence of architectural-scale settlement signatures in the interior itself, not the absence of any interior influence on the coast. Second, exchange across 300--500~km of forested volcanic terrain prior to organised road infrastructure produces a sharply attenuated material-cultural trace compared with exchange across the 8{,}000~km maritime network we document in \S\ref{sec:ch6}; the coastal sites would be expected to look predominantly like West Javanese coastal sites, not like displaced interior sites. Third, the Buni and Batujaya assemblages have not, to our knowledge, been re-examined through the lens of possible interior-origin provenance using modern archaeometric techniques. Doing so is a concrete future research direction the framework predicts would be fruitful. The question is therefore a real one, with a partial answer already in the current record and an open testable component for future work.

This channel is falsifiable. Discovery of open-air pre-400~CE sites in Java's volcanic interior, through systematic survey targeting predicted locations, would weaken or eliminate the within-island contrast. Recovery of organised pre-Hindu Philippine sites at sedimentation rates comparable to Java's 4.4~mm/yr would reduce the cross-regional control. Archaeometric re-examination of West Javanese coastal assemblages that finds no material consistent with volcanic-interior provenance would constrain the interior-to-coastal transfer claim made in the previous paragraph.


\subsection{Linguistic substrate reconstruction}
\label{sec:ch3}

The third channel is linguistic. If a substrate civilisation existed in Java before the historically attested horizon, its material record can be erased but its language cannot be entirely rewritten. What an imposed literary language can do is mask, displace, and relabel the substrate; what it cannot do is delete vocabulary items that the living speech community continues to use. The linguistic channel therefore asks a different question from the preceding two: not ``where are the settlements,'' but ``what does the vocabulary remember?''

The channel sits inside a theoretical frame developed over three decades of Southeast Asian historiography by \citet{wolters1999}, who argued that the Austronesian societies of island and mainland Southeast Asia did not receive Indic, Sinic, or Islamic cultural elements passively but \textit{localised} them, filtering imported material through pre-existing ``place-specific socio-cultural and linguistic coordinates.'' Wolters's framework presupposes a substrate; this channel asks what the substrate looks like when examined directly.

Four independent subquestions return converging answers. First: does Austronesian have a pre-Indic vocabulary for writing and record-keeping? Reconstruction of the Proto-Austronesian lexicon identifies \textit{*surat} ``to write, mark, design'' as an indigenous item at a time depth of approximately 5{,}000~BP, three millennia before any documented Indic contact with Java \citep{blust2009acd}. Related cognates \textit{*tulis} ``to write, paint, decorate'' and \textit{*buku} ``knot, joint'' (later generalising to ``book'') follow the same pattern: the semantic field around inscription and marking has indigenous foundations that predate the Brahmi-derived \textit{aksara} scripts used in the historical period (E112, E113). Whatever was written in those scripts from the 4th century~CE onward was written into a language that already had a vocabulary for writing.

Second: what do kakawin texts---Old Javanese poems composed in court settings between the 9th and 15th centuries~CE and maximally exposed to Sanskrit prestige---show about domain-level vocabulary source? A curated lexicon of 189~kakawin terms examined in earlier work (E058) showed a strong domain gradient: approximately 91\% of agricultural and material-culture vocabulary was indigenous, while approximately 86\% of religious and cosmological vocabulary was Sanskrit. This result, if taken at face value, would be a remarkable signal. We treat it with more caution here, for two reasons disclosed honestly.

A follow-up computational pass applied the same etymological heuristic at dictionary scale, using the Old Javanese Wordnet (5{,}019 synsets; E208). The extremes dampen at this scale: material-culture vocabulary remains majority-indigenous but not at the 91\% level, and religious vocabulary shows more non-Sanskrit native residue than the curated kakawin set suggested. Three explanations contribute. The heuristic used in E208 undercounts Sanskrit loans that have undergone sufficient phonological nativisation to pass as indigenous. The curated 189-term kakawin set is weighted toward literary registers where Sanskrit prestige is strongest. And the kakawin set measures token-level usage in a specific genre, while the Wordnet measures type-level lexical coverage across the full language. The 91\%/14\% figures describe how agricultural and religious vocabulary behave in elite literary production; the 5{,}019-synset scale describes how they behave across the language as a whole. Both are real, and they are consistent with a model in which Sanskrit was genre-selectively adopted by the highest-status literary register while the everyday vocabulary of agriculture, maritime work, domestic technology, and local cosmology remained largely indigenous. Section~\ref{sec:ch3} in a supplementary analysis documents the two scales in full.

Third: do substrate items persist in the modern speech communities of eastern Indonesia? A comparative analysis of six Sulawesi languages identified 438 vocabulary items that do not reduce to Austronesian or Indic sources and that pattern geographically in a way consistent with a pre-Austronesian substrate (E027). The substrate is most dense in agricultural, ritual, and landscape vocabulary---the domains with the deepest continuity between everyday practice and oral tradition. XGBoost classification distinguishes substrate from non-substrate items with AUC~0.76, suggesting that the signal is statistically coherent rather than a product of unprincipled residual labelling. The 230 ``ghost words'' identified in Old Javanese inscriptions as pre-Indic vocabulary that disappears from the written record after the 9th century~CE (E165) reflect the complementary pattern: substrate items that the literary tradition absorbs, transforms, or drops as inscriptional genre conventions harden.

Fourth: can phonological change trace the displacement of substrate material culture by later imports? The PAN \textit{*sagu} ``sago'' reconstructs to Proto-Austronesian as the staple starch of the ancestral speech community. In Old Javanese and its modern descendants \textit{sego}/\textit{sangu} (``rice,'' also ``food generally'') has replaced it, with the phonological evolution \textit{*sagu} $>$ Old Javanese \textit{sagu} $>$ modern Javanese \textit{sego} explained by regular sound change (E198). The substitution tracks the ecological transition from sago-palm economies in the archipelago's eastern edge to rice-agricultural economies in the volcanic-soil west. The linguistic artifact of that transition remains visible in the derivation even where the original ecological system has been lost.

Taken together, the four subquestions converge on a claim the Sanskrit cosmopolis framing of early Indonesia cannot accommodate without modification. Indic vocabulary did not arrive in a linguistic vacuum. It arrived in a speech community that already had words for writing, marking, agriculture, landscape, ritual, and cosmology; that localised the imported vocabulary into a domain-graded pattern visible in the text record; and that preserves substrate items into the modern languages in quantities large enough to reconstruct with conventional comparative methods.

What the channel cannot resolve is as important as what it can. Linguistic depth does not establish material settlement density. A civilisation can have deep vocabulary and thin archaeology if the settlement forms are organic, as the preceding channels argue. And the substrate vocabulary reconstructions we draw on here were assembled by Austronesian linguists for reasons having nothing to do with taphonomy or with VOLCARCH's thesis; their independence from our argument is therefore considerable, but their evidential grain is lexical, not geographic or quantitative.

This channel is falsifiable. A systematic NLP reanalysis of the full kakawin corpus that finds Sanskrit lexical dominance across \textit{all} domains, including agriculture and material culture, would collapse the domain-gradient result. Reconstruction of \textit{*surat}, \textit{*tulis}, and related writing-domain PAN vocabulary as post-Indic loans would remove the pre-Indic writing-concept argument. Documented Austronesian cognates for the substrate items in Sulawesi, traceable to a Proto-Malayo-Polynesian source, would reclassify them from substrate to divergent inheritance.


\subsection{Paleogenomic reinterpretation}
\label{sec:ch4}

The fourth channel is genomic. The primary signal it contributes to this synthesis is positive, not absence-based: modern Javanese whole-genome data shows a depth of lineage-specific variation and substrate ancestry fraction that is incompatible with a recent-migration-only model for population history and therefore requires a substantial, continuously inhabited local population well before the 4th~century~CE.

\citet{maulana2024} sequenced 227 West Javanese whole genomes and identified approximately 1.8~million novel single-nucleotide variants not present in global reference databases---a degree of lineage-specific accumulation that implies either long local habitation or a highly structured local population over a large effective size, or both. Archipelago-wide admixture surveys \citep{lipson2018,larena2021} distribute modern Indonesian ancestry across multiple migratory waves over the past 50{,}000 years rather than a single recent Austronesian replacement, and earlier reconstructions of Island Southeast Asian population history identify substrate non-Austronesian ancestry in western Indonesia at a fraction that requires pre-Austronesian population contribution \citep{lipson2014}. Whatever the precise demographic history, the pre-4th-century-CE population of western Indonesia was not empty. The genetic continuity does not by itself establish the 1--2~million figure from Section~\ref{sec:puzzle}, but it removes one alternative interpretation of that section's gap: the Javanese population we are arguing for is not hypothetical in its genetic traces; it is the population modern Javanese descend from.

The channel's second contribution concerns the pattern of ancient DNA \textit{preservation} across Southeast Asian contexts. Here we flag a tautology the reader may reasonably notice before we do. To our knowledge, no Holocene or historic-period human aDNA has been published from any open-air archaeological context in Java's volcanic interior. But the reason is partly that no securely dated pre-400~CE human skeletal remains have been identified in such contexts to test. One cannot recover aDNA from bones one does not have. The absence of aDNA is therefore not independent evidence of a preservation filter acting on DNA alone; it is a downstream consequence of the same broader taphonomic regime that leaves no skeletal material above the pre-Hindu horizon for any analysis.

The informative comparison is not within-Java but between Java and Southeast Asian regions with contrasting preservation contexts. Ancient DNA has been recovered from Wallacean karst-cave contexts---the canonical case is the Leang Panninge individual from South Sulawesi, sequenced by \citet{carlhoff2021} from petrous bone in a limestone-cave burial dated to approximately 7{,}200--7{,}300~years~BP---and from Liang Bua (Flores), Hoabinhian sites on mainland Southeast Asia \citep{mccoll2018,lipson2018}, and Philippine cave sites \citep{larena2021}. What this cross-regional comparison shows is that Southeast Asian tropical climate is not an absolute barrier to aDNA recovery when preservation contexts are favourable, and that every successful Southeast Asian aDNA recovery to date has come from karst-cave rather than open-air tephra-mantled contexts. Java's volcanic interior provides few karst-cave microenvironments and extensive tephra-mantled open-air deposits; the cross-regional pattern is consistent with what the compound-taphonomy framework would predict about where skeletal remains (not only DNA) should survive.

The three observations together---(i) modern West Javanese WGS shows deep lineage-specific variation, (ii) archipelago-wide admixture analyses require multiple pre-Austronesian waves of settlement in western Indonesia, (iii) Southeast Asian aDNA recovery success pattern maps onto preservation-context availability in a way that predicts Java's silence---constitute the genomic channel's contribution. The first two are positive signals; the third is a consistency check on the framework. None of the three is, by itself, proof of the specific population magnitude argued elsewhere in this paper. Together they remove the alternative interpretation that modern Javanese are recent arrivals inhabiting a previously empty volcanic landscape.

The channel's principal independence from the preceding three is that the underlying data was produced by paleogeneticists whose research interests are Austronesian dispersal, Wallacean hunter-gatherer ancestry, and modern Southeast Asian admixture history. None of these research programmes is concerned with Javanese taphonomy. The reinterpretation here takes published data at face value and asks what deep ancestry and the regional preservation pattern together imply about the plausibility of a substantial pre-400~CE population, not about the specific magnitude or form that population took.

What the channel cannot resolve is the population magnitude. Genetic data place lower bounds on lineage continuity; they do not estimate absolute headcount. The 1--2~million central estimate from Section~\ref{sec:puzzle} is not a genetic result and cannot be derived from the genomic record. The channel's role is eliminative: it removes the alternative interpretation under which modern Javanese are recent arrivals inhabiting a landscape that was previously empty. It does not, on its own, discriminate between a population of 300{,}000 and one of 3~million.

This channel is falsifiable. Evidence that modern Javanese WGS lineage-specific variation traces principally to post-1st-millennium-CE migrations rather than deeper continuity would collapse the primary positive signal. A re-analysis of \citet{lipson2014} or \citet{larena2021} showing that substrate non-Austronesian ancestry in western Indonesia can be accounted for entirely by post-Austronesian-expansion admixture sources would remove the second. Systematic publication of East Javanese whole-genome data showing higher lineage diversity and older admixture than Maulana et al.'s West Javanese sample would reverse the framework-based prediction that eastern volcanic-interior populations experienced more taphonomic bottlenecking.


\subsection{Colonial archive text mining}
\label{sec:ch5}

The fifth channel mines the archaeological reports and newspaper coverage of the Dutch colonial period for material that the original observers could not themselves interpret in the framework we apply here. Between 1854 and 1949 Dutch fieldworkers, colonial administrators, and Javanese informants produced a sustained textual record of archaeological encounters in Java: the \textit{Oudheidkundig Verslag} (OV) annual reports of the colonial antiquities service, the \textit{Delpher} newspaper archive of the Royal Library of the Netherlands, and associated publications. These texts describe finds, locations, conditions, depths, and circumstances in a vocabulary appropriate to their own moment: artefacts are catalogued, sites are mapped, and occasional depth measurements are recorded when construction work exposed them. The observers were not thinking about cumulative volcanic sedimentation, nor were they adjusting their spatial sampling to control for burial-depth gradients. Their reports are therefore a pre-VOLCARCH, pre-computational, human-curated corpus that can be examined post-hoc for the patterns the framework predicts.

Two operational extractions were performed. A named-entity-recognition pass over 16~OV volumes (1925--1949) produced 22{,}162~candidate settlement-mention tokens, of which a subset with sufficient contextual signal have been linked to modern geographic coordinates (E091). A broader Delpher query yielded 1{,}768~newspaper articles between 1854 and 1942 mentioning archaeological finds in Java, of which 165 have so far been geocoded and 33 have sufficient textual specificity to reconstruct burial depth (E141). The NLP methodology required extensive handling of pre-1947 Dutch orthography and continues under development for the colonial historical register \citep{verberne2024}.

Three patterns emerge from this corpus that have not been explicitly reported in the primary colonial literature. First, a volcano-distance gradient in find frequency: of geocoded colonial-era finds, the density in the 0--15~km radius around active volcanic centres is approximately 2.4\% of the total---substantially below the value expected under uniform geographic distribution and consistent with the VOLCARCH prediction that the immediate volcanic flanks should be archaeologically suppressed even in retrospective record compilation. Second, a depth distribution: colonial depth measurements, where recorded, cluster at 1.0--4.0~m with a median near 2.5~m---matching the depth that a 4~mm/yr cumulative sedimentation rate would produce over a 760-year projection (Section~\ref{sec:ch1}). The colonial observers, documenting finds exposed by building construction, irrigation works, or erosion, did not design this distribution; the match with the sedimentation calibration is unintended convergence. Third, a spatial enrichment: the geocoded colonial finds are approximately 5.8 times more likely to lie within VOLCARCH-predicted buried-settlement zones than chance would predict (E197).

The channel's independence is structural rather than methodological. The corpus was produced 77~to 172~years before the present analysis, by observers with their own interpretive framework and their own biases (broadly Dutch-colonial archaeological priorities, with emphasis on monumental Hindu-Buddhist sites over settlement archaeology), using no computational tools of any kind. The mining here does not modify the texts; it re-reads them. The biases of the corpus are visible and can be partially corrected. The biases of the mining itself---the choice to look for sedimentation-rate signatures and volcano-distance gradients---are VOLCARCH-endogenous and are flagged honestly: we find what we looked for.

What the colonial archive cannot resolve is the question of whether the finds it documents are the tip of a buried settlement matrix or the full surface-exposure record of a smaller substrate. Colonial depth records are biased toward the shallower end: construction projects rarely went deeper than 2--3~m before the mid-20th century, and finds deeper than that appear only in the few instances where irrigation or foundation work exceeded that threshold. The 1.0--4.0~m range reported here is therefore an observation-window artefact, not a real upper bound on pre-Hindu deposit depths.

This channel is falsifiable. A systematic re-geocoding of the 22{,}000+~OV mentions that reveals no volcano-distance gradient and no depth-calibration match would remove the spatial signal. A reanalysis using a corpus without VOLCARCH-informed query design---an independent NLP group applying generic archaeological extraction to the same OV volumes---that fails to reproduce the 2.4\% near-volcano density, the 2.5~m median depth, or the 5.8$\times$ enrichment would collapse the empirical basis for the channel's conclusions.


\subsection{Archaeometric evidence from exported Javanese material culture}
\label{sec:ch6}

The sixth channel is new to this synthesis, added in response to a systematic literature review of the archaeometric, bronze-drum, and trans-Eurasian trade literature conducted independently of the VOLCARCH experiment programme. The first five channels examine what is \textit{not} there in volcanic Java. This channel examines what \textit{is} there---in Bali, China, Korea, Egypt, and Micronesia, in tombs and ports spanning the 1st through 6th centuries~CE, produced in Java and Bali and exported along the maritime networks that connected those regions.

The channel has two sub-channels, each with its own specialist research community and each addressing a distinct phase of the invisibility regime argued in this paper: Pejeng-type bronze drums, which date to approximately the 1st--3rd centuries~CE and therefore fall squarely in the pre-inscription period, and \textit{Jatim} glass beads, whose documented extra-Javanese distribution dates to approximately the 4th--6th centuries~CE and therefore overlaps the emergence of Tarumanagara and the early Hindu-Buddhist polities. We separate these chronologically because they support different parts of the argument; collapsing them would overclaim.

\paragraph{Sub-channel 6A: Pejeng-type bronze drums of local Bali/Java production (pre-inscription).} Bronze drums of the Pejeng type---stylistically distinct from mainland Dong Son through their hourglass body form and two-piece mantle-plus-tympanum casting---are distributed across Bali, Java, and eastern Indonesia \citep{bernetkempers1988}. \citet{calo2014} applied lead-isotope analysis to the Manikliyu drum (Pejeng typology, Bali) and demonstrated that its isotopic signature is consistent with raw-material sourcing from mainland Southeast Asian Dong Son networks while its morphology and casting technique are local Bali/Java innovations. Subsequent work has extended compatible analyses to Timor-Leste finds \citep{calo2019}. The largest surviving instance, the ``Moon of Pejeng,'' at 186.5~cm in diameter, is the largest bronze drum documented in Southeast Asia and cannot have been produced by a non-industrial workshop.

The dating window for Pejeng production is commonly placed in the late Bronze Age to early Iron Age Southeast Asian sequence, approximately the 1st through 3rd centuries~CE on typological and regional-network grounds \citep{bernetkempers1988,calo2014}. This window is pre-inscription in Java: it precedes Tarumanagara (mid-4th century~CE) and Kutai (approximately 400~CE). For this channel's purposes, the Pejeng drums establish that an archaeometrically traceable, technologically organised production tradition existed in Bali and Java before any documentary archaeological horizon and during a period for which volcanic Java's interior shows no settlement record whatever. The workshops themselves are archaeologically absent, but their outputs, in the specific form of large cast objects whose lead-isotope signatures identify their raw material provenance and whose morphology identifies their production tradition, survived the same taphonomic processes that appear to have erased the settlement matrix.

\paragraph{Sub-channel 6B: \textit{Jatim} glass beads in elite extra-Javanese tombs (invisibility through the early state period).} \citet{jia_etal_2024} reported non-destructive compositional analysis of two glass eye-beads excavated from the Dongxin tomb complex at Datong, northern China, dated to the Northern Wei Pingcheng period (398--494~CE). The beads are chemically identifiable as \textit{Jatim} products. The authors' diagnostic rests on the simultaneous presence of four features: a mixed plant-ash soda-lime (v-Na-Ca) base for the matrix; mineral-soda alumina (m-Na-Al) high-alumina regions in the eye decorations; millefiori eye-bead morphology specific to the \textit{Jatim} typology \citep{lankton2008}; and a colorant palette drawing on cobalt with low manganese (indicating Mediterranean sourcing rather than Chinese sourcing), SnO\textsubscript{2} tin-white opacifier, Cu\textsubscript{2}O red, and PbSnO\textsubscript{3} lead-tin yellow---all technologies either absent or incipient in 5th-century Chinese glass production. We note the methodological caveat explicitly: neither v-Na-Ca nor m-Na-Al alone is diagnostic of Java \citep{jia_etal_2025,wang_etal_2023}. The Java attribution depends on the full signature, not on any single feature, and any future application of the method to new finds must maintain this discipline.

The Datong find is not an isolated anecdote. \citet{lankton_bernbaum} document at least ten \textit{Jatim} beads from elite tombs in Gyeongju, the Silla Korean capital, dating from the late 4th to the mid-6th century~CE. One of these---the bead comprising Korean National Treasure 634 from the Sikrichong tomb---has received full LA-ICP-MS, SEM-EDS, and EPMA characterisation and provides the canonical archaeometric reference for subsequent \textit{Jatim} identifications. \citet{then2014} reports \textit{Jatim}-type and Indo-Pacific beads from the Late Harbor Temple deposits at Berenike, the Roman Red Sea port excavated by \citet{sidebotham2011}, dated to the 4th through early 6th century~CE. Further distribution points, with typological rather than archaeometric confirmation, include Palau in Micronesia (where \textit{Jatim} beads circulated as currency) and sites across the western Roman world \citep{green2018}.

We frame this sub-channel's contribution carefully. The \textit{Jatim} distribution dates primarily to the 4th through 6th centuries~CE, which overlaps the emergence of Tarumanagara in coastal West Java and the earliest Hindu-Buddhist polities. It therefore does not directly demonstrate the existence of a pre-Hindu civilisation; reviewers of an earlier version of this paper correctly noted that using the \textit{Jatim} corpus as a pre-Hindu anchor would be a chronological sleight of hand. What the \textit{Jatim} evidence does demonstrate is that the invisibility regime documented by Channels 1--5 for the pre-400~CE horizon \textit{extends through} the early state-formation period, in the specific sense that the production workshops that sustained a trans-Eurasian glass export system between approximately 400 and 600~CE are themselves archaeologically absent from Java and Bali. The same taphonomic filter that erases pre-Hindu settlements does not suddenly become permeable when Tarumanagara's first stone inscriptions are erected; it continues to erase organic workshop infrastructure even as monumental stone announces a new regime. The \textit{Jatim} sub-channel thus strengthens Channel~1 (the burial process is continuous) and adds empirical content to the inference that the pre-Hindu interior, under a similar filter, would also have supported organised production that left no surface-recoverable trace.

\paragraph{Channel independence, convergence, and scope.} Considered jointly, the two sub-channels constitute a distinct form of evidence whose independence from VOLCARCH's evidential base is structural. The Manikliyu lead-isotope analysis was conducted by metallurgists working through \citeauthor{bernetkempers1988}'s mid-20th-century typological foundation. The Datong excavation was conducted by Chinese archaeologists for reasons internal to Northern Wei period studies. The Gyeongju corpus has been studied by Korean archaeologists and by \citeauthor{lankton2008}'s canonical Southeast Asian bead-typology programme. Berenike is excavated by a University of Delaware team studying Roman Red Sea trade. None of these research programmes has any stake in the VOLCARCH thesis. Their finds speak to their own questions. What this channel adds is that the invisibility regime---organised production present in extra-Javanese distribution but archaeologically absent in Java/Bali itself---holds across two object classes, two specialist communities, and at least six centuries, from Pejeng's pre-inscription dating into the post-Tarumanagara Jatim horizon.

The sub-channels also differ in what they claim. Sub-channel 6A (bronze) is VOLCARCH's pre-Hindu anchor within this channel: it documents a 1st--3rd-century~CE production tradition in a period for which Java's archaeological interior is empty. Sub-channel 6B (glass) does not extend that anchor further back; rather, it extends the invisibility-of-workshops pattern forward into the early state period, demonstrating that the phenomenon Channels 1--5 argue for is not confined to a pre-Hindu ``darkness'' that lifts when inscriptions appear.

What the channel cannot resolve is the location or form of the production centres in either period. Archaeometry establishes origin and chronology but cannot substitute for excavation of the workshops. The framework here is also vulnerable to the attribution caveats surfaced in the review literature: if future studies demonstrate that the mixed v-Na-Ca + m-Na-Al + Mediterranean-colorant signature can be produced in South Asian or mainland Southeast Asian workshops at comparable specificity, the Java attribution of individual Jatim beads and therefore the chronology of Javanese glass production would require revision. Current specialist consensus treats the Java attribution as secure \citep{lankton2008,jia_etal_2024}; the channel's strength depends on that consensus holding.

This channel is falsifiable. A lead-isotope or trace-element reanalysis of the Pejeng drums that identifies a non-Bali/Java workshop origin would collapse Sub-channel 6A and remove the pre-Hindu anchor. An alternative production origin for the Datong, Gyeongju, and Berenike \textit{Jatim} beads that survives peer review at the methodological standard of Jia et al.\ (2024) would collapse Sub-channel 6B. Re-dating of the Pejeng drum tradition to post-4th century~CE, on stratigraphic grounds not yet published, would remove its status as pre-inscription evidence and collapse the chronological load the sub-channel is asked to bear.


\subsection{Convergence across channels}
\label{sec:convergence}

Six channels, each with its own data, methods, and observer community, converge on a single pattern. The sedimentation calibration establishes that pre-400~CE deposits sit at 4--10~m in volcanic Java, below surface-survey reach. The Philippines comparison shows that zero documented pre-400~CE sites in volcanic Java contrast with 275--340 in a neighbouring Austronesian region with the same organic-architecture tradition but lower volcanic burial and more karst preservation. The linguistic substrate channel shows indigenous vocabulary with pre-Indic depth, domain-graded patterning in the literary record, and substantial substrate persistence in modern speech communities. The genomic channel provides the positive signal of deep lineage-specific variation in modern Javanese WGS and admixture-depth patterns incompatible with recent-migration-only ancestry, and documents a Southeast Asia-wide preservation pattern (aDNA success in karst, failure elsewhere) consistent with the taphonomic framework. The colonial archive channel surfaces, in pre-computational 19th- and 20th-century Dutch records, a volcano-distance gradient and a depth distribution matching the sedimentation calibration. The archaeometric channel documents, in two chronological layers, first a 1st--3rd-century~CE local Pejeng bronze drum production tradition (the pre-inscription anchor, Sub-channel 6A) and second a 4th--6th-century~CE Javanese \textit{Jatim} glass bead export corpus (showing the invisibility regime extends through the early state period, Sub-channel 6B).

No single channel is decisive. Sedimentation alone cannot distinguish buried-and-absent from buried-and-present. Philippines comparison alone cannot control for every ecological and institutional confound. Linguistic depth alone does not map to geographic scale. Paleogenomic reinterpretation alone cannot establish absolute population. Colonial archive mining alone is vulnerable to corpus-specific biases and observation-window artefacts. Archaeometric evidence alone could in principle be explained by workshops outside Java that were misattributed, or by a small production enclave not implying broader civilisation.

What the six channels accomplish together is the narrowing of the space of alternative explanations to a point where a purely developmental reading of the Javanese archaeological record---Hindu-Buddhist kingdoms crystallising from a quantitatively sparse pre-existing population---is quantitatively incompatible with the joint signatures. The pattern that fits all six channels simultaneously is the one this paper develops in the remaining sections: selective survival of durable material against compound taphonomic erasure of the organic settlement matrix, in a substrate civilisation whose population was on the order of one to two million at 400~CE and whose production systems left traces in geographies where the taphonomic filter did not operate.


%% ================================================================
\section{Selective Survival}
\label{sec:selective_survival}

The six channels assembled in Section~\ref{sec:channels} establish a pattern: organised production traces survive where the taphonomic filter does not operate (extra-Javanese tomb contexts, karst-preserved aDNA, modern-day lexical substrate), while the settlement matrix that produced them survives nowhere in Java's volcanic interior. The pattern is not one of uniform erasure. It is selective, and what survives is informative.

The concept organising this section is \textit{selective survival}. An archaeological record is the joint product of what was produced and what the preservation filter let through. When the filter is severe and materially selective---as it is in volcanic tropical contexts with an organic architectural tradition---the resulting record is not a scaled-down version of the original material culture but a qualitatively biased sample of it. Specifically: durable materials (bronze, stone, fired ceramic in protected contexts, glass in elite tombs far from the production zone) survive at rates approaching unity across centuries; semi-durable materials (iron, unfired or unprotected ceramic, shell) survive at modest rates with heavy condition loss; organic materials (wood, bamboo, palm thatch, plant-based textiles, food remains, organic pigments) survive at effectively zero rates in open-air tropical contexts. A society whose architecture, household inventory, subsistence record, and most of its technological expression were organic will therefore appear, to archaeology alone, as a society whose every material output was metal or stone. What remains in the ground is not the civilisation; it is the civilisation's shadow on the only materials durable enough to leave one.

The bronze drum record illustrates the point. Approximately forty nekara of documented Javanese provenance and thirty or more from eastern volcanic territories, including the Tuban specimen (Heger Type II, approximately 300~BCE, containing a small bronze elephant and axe possibly of local manufacture) and the Manikliyu Pejeng drum discussed in Section~\ref{sec:ch6}, represent the Bronze Age material record of Java and Bali that has survived to the present. These are not small objects. The Moon of Pejeng at 186.5~cm diameter could not have been cast by a non-industrial workshop. The forty surviving nekara imply---under any reasonable estimate of bronze-to-other-materials ratio in an organic-architecture society---a material production system on the order of thousands of non-durable objects for each surviving drum. The drums are not the record; they are the 1\% of the record that the filter passed.

Table~\ref{tab:survival_matrix} summarises the pattern by material class, drawing on standard tropical-taphonomy literature \citep{erasmus2005}.

\begin{table}[t]
\caption{Material-class survival through the compound taphonomic filter operating in volcanic tropical Java. Survival rates reflect open-air contexts; karst-cave contexts preserve substantially more of the semi-durable and organic classes but, as established in Section~\ref{sec:ch2}, are minimally available in Java's volcanic interior.}
\label{tab:survival_matrix}
\begin{tabular}{p{3.5cm}p{2cm}p{2cm}p{5.5cm}}
\toprule
Material class & Tropical decay & Volcanic burial & Archaeological consequence \\
\midrule
Bronze & resistant & resistant & survives at approximately 100\% over millennia in buried contexts; \textasciitilde40 nekara recovered from Java/Bali \\
Stone (andesite, limestone) & resistant & resistant (partial re-exposure possible) & monumental survivors; Dwarapala statues, megalithic features \\
Fired ceramic & semi-resistant & fragments survive; crushed & sherds recoverable, whole vessels rare \\
Iron & oxidises rapidly & partial survival & corroded fragments; hoe-blades, knife-fittings \\
Wood & destroyed & destroyed & zero survival in open-air tropical contexts \\
Bamboo / palm thatch & destroyed & destroyed & zero survival; entire architectural register invisible \\
Plant textiles & destroyed & destroyed & zero survival \\
Organic residues, foodways & destroyed & destroyed & zero survival outside protected contexts (phytoliths possible) \\
\bottomrule
\end{tabular}
\end{table}

The selective-survival framing reframes the central argument of this paper in a way that is both stronger and more honest than the conventional absence-of-evidence formulation. We do not claim that volcanic Java has no pre-Hindu archaeology. Volcanic Java has a pre-Hindu archaeology consisting predominantly of bronze drums (Channel 6 Sub-channel 6A), some fragmentary fired ceramic, and a few stone megalithic features, all of which are visible in the existing record. What volcanic Java lacks---and what demographic and comparative considerations require---is the \textit{organic settlement matrix}: the wooden and bamboo dwellings, the palm-thatch villages, the wet-rice field systems, the stored foodstuffs, the plant-based textiles, the ritual implements in non-durable materials. These are the materials that survive nowhere in open-air tropical contexts, and the taphonomic filter in volcanic Java is additionally critical to even the semi-durable classes.

The framing has an implication for what further evidence could look like. If the framework is correct, systematic deep coring or construction-triggered deep excavation in predicted burial zones should encounter, at appropriate stratigraphic depths, precisely those semi-durable and occasionally durable artifacts that remain from an eroded organic society: pottery fragments, charcoal, ash lenses, bronze fragments, iron corrosion products, and, in favourable microenvironments, phytoliths. These are modest expectations by the standards of a freshly excavated site, but they would constitute a qualitative shift from the current pre-Hindu record. The falsification criterion attached to this prediction (Section~\ref{sec:predictions}) specifies its operationalisation.

We note what the framing does \textit{not} claim. It does not predict that pre-Hindu Java produced monumental architecture on a scale comparable to the Hindu-Buddhist period. The absence of stone temples before the 4th century is likely real, not a preservation artifact: stone survives, and none of a pre-Hindu monumental scale has been recovered. The claim is specifically about the organic settlement matrix of a population supported by volcanic-soil agriculture, not about the institutional or monumental expression of that population. The substrate civilisation we argue for is one of villages and workshops, not of pre-Hindu kingdoms. Our argument on material production (Section~\ref{sec:ch6}) concerns technologically organised craft output; it does not require, and we do not claim, pre-Hindu urbanism or state-level hierarchy.

One caveat requires explicit acknowledgement before proceeding. The selective-survival framing could, without discipline, become a container that absorbs any outcome: if sites are found, the framework is vindicated; if no sites are found, the filter was even more severe than estimated. To avoid this, we specify ahead of the predictions section (\S\ref{sec:predictions}) the quantitative thresholds that would constitute falsification rather than rescue of the framework. The framing itself is not a hypothesis; the specific predictions and their pre-registered thresholds are.


%% ================================================================
\section{Living Evidence: Wayang as an Archaeological Record in Motion}
\label{sec:wayang}

Most archaeology is dead. Wayang is performed nightly. Its characters, cosmological structure, and narrative repertoire encode cultural forms whose material correlates are invisible in the ground, but whose continuity in the living tradition is documented in centuries of performance practice. We include wayang in this synthesis not as folklore illustration but as a form of evidence that speaks to the same question the preceding channels address: what was there before the Hindu-Buddhist horizon, and how did it transform rather than disappear?

The wayang repertoire divides into three layers with sharply different provenances \citep[E205]{amien_e205}. The \textit{pakem} (trunk) comprises the Mahabharata- and Ramayana-derived narrative core that arrived with Indic contact and represents approximately 30--40\% of current performance content. The \textit{carangan} (branch) comprises Javanese elaborations on Indic frames---plotlines that insert Javanese characters and settings into Indian narrative structures---at approximately 30--40\%. The \textit{sempalan} (twig) comprises stories with no Indian source: entirely Javanese in character inventory, setting, and narrative logic, at approximately 20--30\% of the active repertoire. Sempalan stories are not marginal curiosities. They include some of the most performed and ritually significant lakon in Central and East Javanese puppet theatre. Their existence at 20--30\% of a repertoire dominated on the surface by Indic content is direct evidence for a substantive pre-Indic narrative tradition that the Hindu-Buddhist cultural overlay absorbed rather than replaced.

The Punakawan (clown-servant) character complex is the sharpest instance. Semar, Gareng, Petruk, and Bagong in the Central and East Javanese tradition, with Togog and Bilung in some Central Javanese variants, form a quartet or quintet of characters whose etymology, cosmological role, and performance function are all Javanese rather than Indic. Semar's name traces to a Javanese root, not a Sanskrit one; his cosmological role is explicitly pre-Hindu (a deity demoted to servant form by the imposition of Indic narrative structures); his performance function is as ritual commentator and moral authority, not as comic relief. In the 1358~CE relief programme of Candi Tigamangi in East Java, Semar appears in a form consistent with the modern wayang character, predating the codification of the Surakarta puppet tradition by four centuries. Scholarly consensus treats Semar as a pre-Hindu Javanese deity whose Indic-period fate was ritual demotion rather than elimination; he retains, in the living tradition, a sacred authority that the hero-prince protagonists formally senior to him do not \citep[see the extensive treatment in][]{brandon1970}.

Parallel characters in Balinese (Twalen), Kelantanese Malay (Pak Dogol and Wak Long), and Sundanese (Cepot) wayang traditions reinforce the pattern across the Austronesian cultural sphere. The Punakawan complex is not a local Javanese invention; it is a shared Austronesian substrate feature expressed through locally distinct puppet-theatre conventions.

The \textit{gunungan}, the mountain-tree puppet that opens and closes every performance, carries comparable archaeological content in compressed symbolic form. Its visual vocabulary includes elements drawn from Austronesian cosmology (the world mountain, the banyan tree, the ocean below, the local banteng and other native fauna) and frames them with Indic architectural motifs (\textit{kala-makara} heads above and around the central field). The gunungan is both a material artifact and a semantic object: its iconographic programme encodes the same cultural layering that the wayang repertoire encodes narratively and that Javanese language encodes lexically. Living cultural objects are not normally treated as archaeological evidence; the argument for including the gunungan here is that its iconographic elements are systematic and stable enough to constitute a text legible through conventional comparative methods, and that the text's content is precisely the pre-Hindu cosmological material that the archaeological record of volcanic Java does not preserve \citep[see][]{kieven2013}.

Wayang shows the reader what Java's volcanic interior archaeology could look like if organic culture survived. The puppets themselves are materially rich: buffalo-hide figures painted with natural pigments, carved-and-incised decoration, gold leaf on royal puppets, gamelan bronze accompaniment, textile staging, elaborate ritual food offerings---an entire material culture whose ground-surface equivalent for pre-Hindu centuries has been erased. The performance tradition documents non-Indic ritual specialists (the \textit{dalang}), indigenous musical theory in the gamelan tuning systems (\textit{sléndro} and \textit{pélog}, whose structural logic is non-Indic and reconstructible from comparative Austronesian metrical and scalar systems), and geographic and technological knowledge in the lakon texts (place names, travel sequences, material culture mentions) that a future NLP programme could extract systematically from the published lakon corpus. We note this extension as a concrete future research direction consistent with the methodological approach of this paper.

The argument we draw from wayang is modest and specific. The living tradition is not proof of a pre-Hindu civilisation; it is evidence that pre-Hindu cultural content persisted in sufficient volume and structure to constitute a measurable residue within a post-Hindu performance tradition four to six centuries after its originating conditions. That residue would not exist if the originating conditions had been demographically trivial. Twenty to thirty percent of an active wayang repertoire is a lot of story to preserve across more than a millennium of cultural layering; the channel complements, rather than substitutes for, the material-culture and linguistic channels developed earlier.


%% ================================================================
\section{A Six-Filter Framework of Archaeological Erasure}
\label{sec:layers}

The six channels of Section~\ref{sec:channels} and the selective-survival framing of Section~\ref{sec:selective_survival} can be organised into a single framework of erasure filters that together produce the archaeological record visible in the present. We present the framework as six filters, acknowledging a seventh confounding factor (karst availability) that does not operate as a filter in the strict sense but modifies the relative intensity of the others. Each filter is documented elsewhere in this paper or in cited literature; their joint operation is what this section formalises.

The filters are: \textbf{L1 Volcanic burial}, which removes open-air organic settlement matter and semi-durable artifacts from surface visibility across the volcanic interior (Section~\ref{sec:ch1}); \textbf{L2 Coastal submersion}, which has placed portions of the pre-Holocene Sunda Shelf under present-day sea level, erasing the coastal settlement record of communities that predated the sea-level rise \citep{oppenheimer1998}; \textbf{L3 Historiographic bias}, under which the archaeological literature has prioritised named kingdoms with monumental and epigraphic records over unnamed polities and oral-tradition communities; \textbf{L4 Cosmological overwrite}, under which Sanskrit and Sanskrit-derived cultural-religious vocabulary and practice have displaced indigenous ritual forms in the literary and institutional record without entirely eliminating them in the living tradition (Section~\ref{sec:wayang}, and the substrate vocabulary evidence of Section~\ref{sec:ch3}); \textbf{L5 Genre taphonomy}, under which the documentary record is dominated by royal-grant and formulaic inscriptional text types that systematically exclude material-culture, subsistence, and daily-life reference; and \textbf{L6 Historiographic periodicity}, under which the accepted chronology of Southeast Asian cultural development has been structured around events visible in the surviving record (Indianisation, monument construction, dynastic cycles) rather than around the underlying demographic and productive systems that those events punctuate.

Table~\ref{tab:filters} summarises the filters, their mechanisms, and the evidence for each.

\begin{table}[t]
\caption{Six filters of archaeological erasure and the karst-availability confound. ``Passes'' indicates what the filter does not block; ``blocks'' indicates what it erases. The filters operate in parallel; the joint product is the record visible to present-day archaeology.}
\label{tab:filters}
\begin{tabular}{p{1.6cm}p{2.6cm}p{3.2cm}p{3cm}p{3cm}}
\toprule
Filter & Type & Passes & Blocks & Evidence \\
\midrule
L1 Volcanic burial & Physical & Metal, stone, deep-cave deposits & Organic open-air settlement & Section~\ref{sec:ch1}; P1-core \\
L2 Coastal submersion & Physical & Upland & Shelf-continent coastal (2.09M~km\textsuperscript{2}) & \citealp{oppenheimer1998} \\
L3 Historio\-graphic bias & Interpretive & Named kingdoms, inscriptions & Unnamed polities, oral traditions & Kutai precedence; Srivijaya capital uncertainty \\
L4 Cosmological overwrite & Cultural & Sanskritised surface & Indigenous substrate & Section~\ref{sec:ch3}; Section~\ref{sec:wayang} \\
L5 Genre taphonomy & Documentary & Royal grants, formulaic texts & Daily life, material culture, oral knowledge & E048, E057 (DHARMA genre analysis) \\
L6 Historio\-graphic periodicity & Temporal & Indianisation-as-linear narrative & Wave structure, localisation cycles & E030 (pre-Indic vocabulary increases over time) \\
\midrule
\textit{Karst availability} & \textit{Confound} & \textit{Cave-preserved deposits} & \textit{---} & \textit{Section~\ref{sec:ch2}; E178 Philippines} \\
\bottomrule
\end{tabular}
\end{table}

The filters combine multiplicatively in their effect on archaeological visibility, and this multiplicativity is what makes the volcanic-Java case severe. A community passed through L1 alone (volcanic burial) would leave durable material culture survivors (bronze drums, stone monuments) detectable by surface survey with some probability proportional to burial depth. A community passed through L1 and L4 (cosmological overwrite) loses not only its open-air organic matter but the linguistic-cultural frame within which that matter was produced and interpreted. A community passed through L1, L3, L4, and L5 further loses its documentary voice: the surviving inscriptional record will not describe its settlements or daily life, only its elite-ritual-political practices. A community passed through all six filters will be archaeologically invisible in a way that the sum of the filters' effects, rather than any single filter, explains.

A complete cascade model that assigns numerical visibility coefficients to each filter has been developed in earlier VOLCARCH work \citep[E110, E176]{amien_e110} and tested against cross-regional comparisons (E155). We have elected not to present the cascade's specific parameter values here because earlier reviewers have correctly identified that multiple parameterisations of a six-filter multiplicative model can bracket the observed archaeological record, making the cascade diagnostic rather than predictive. The diagnostic use of the cascade---identifying which filter contributes most leverage to visibility reduction, which turns out under most parameterisations to be the combination of L1 (burial) and the conventional survey effort directed at its consequences---is what the framework can support, and we present it in that more modest form here.

The karst-availability confound requires separate treatment. As Section~\ref{sec:ch2} documented, the Philippines comparison between volcanic and non-volcanic regions is complicated by differential karst availability. Philippine volcanic zones have documented pre-400~CE sites because those zones have cave-preserved sites that bypass L1 (and, to varying extents, L4 and L5). Java's volcanic interior has minimal karst. The karst confound does not operate as one of the six filters; it operates as a modulator of L1's severity, converting an L1-active region into an L1-suppressed one. A fuller cross-regional analysis of the framework would treat karst availability as a covariate in the visibility model.

The framework's role in this paper is organisational rather than predictive. It does not replace the channels of Section~\ref{sec:channels}; it names the erasure processes those channels collectively indirect. What the framework claims is that the archaeological invisibility of pre-400~CE volcanic Java is a systematic consequence of filters whose operation is documented independently, and that recovering the record those filters produced will require methods matched to the filter stack rather than surface survey adequate for a less-filtered record.


%% ================================================================
\section{Pre-Registered Falsifiable Predictions}
\label{sec:predictions}

We specify eight predictions whose outcomes can confirm, constrain, or falsify the framework developed in the preceding sections. Each prediction is stated with its method of test, the expected outcome if the framework is correct, the expected outcome if it is wrong, a decision threshold, and a rough cost estimate. We commit to these predictions in their pre-registered form: outcomes that contradict them will be reported as failures of the framework, not rescued by auxiliary hypotheses.

\textbf{Prediction 1: Ground-penetrating radar (GPR) in twenty predicted burial zones.} Target selection based on Channel 1 sedimentation rates plus Channel 2 within-island preservation (coastal-adjacent basins with projected pre-400~CE deposits in the 2--5~m depth range, within GPR effective reach). Expected outcome if framework correct: 2.5 anthropogenic-return finds across 20 sites \citep[0--6 range,][]{amien_e116}, with $P(\mathrm{zero}) \approx 7\%$. Expected outcome if wrong: zero returns across all 20 sites. Cost estimate: USD 40{,}000--100{,}000.

\textbf{Prediction 2: Deep coring at West Java coastal paleo-river outlets.} Five priority locations identified from paleo-river reconstructions \citep[E177]{amien_e177} of Sunda Shelf displacement endpoints. Expected outcome if framework correct: one to two organic deposits with dated material matching the pre-Hindu horizon, plus non-zero pottery or charcoal at projected depths. Expected outcome if wrong: no organic material at projected depths. Cost estimate: USD 30{,}000--60{,}000.

\textbf{Prediction 3: East Javanese whole-genome sequencing.} Sequencing 100--200 individuals from East Javanese volcanic-interior populations (e.g., Malang, Kediri, Pasuruan) to a comparable depth and standard as \citet{maulana2024} for West Java. Expected outcome if framework correct: lineage diversity below the West Javanese baseline, admixture timing shifted toward more recent events (consistent with more severe volcanic-era population bottlenecks). Expected outcome if wrong: equal or greater lineage diversity and older admixture signatures in East Java than West Java. Cost estimate: USD 50{,}000--150{,}000.

\textbf{Prediction 4: Systematic excavation at Buni coastal-to-interior transect.} A north-south survey spanning the Buni coastal complex and proceeding inland through the modern Bekasi--Jakarta corridor to non-volcanic but upland zones. Expected outcome if framework correct: a clear preservation gradient from well-preserved coastal (Buni-like) to partially preserved (transitional) to increasingly erased (approaching volcanic influence), with material-culture continuity detectable across the gradient. Expected outcome if wrong: abrupt or absent gradient, or gradient with reversed polarity. Cost estimate: USD 80{,}000--200{,}000.

\textbf{Prediction 5: Kakawin corpus NLP for material culture and geography.} A systematic NLP extraction of material-culture nouns, geographic place names, and technological references from the published Old Javanese kakawin corpus (Zoetmulder and subsequent editions). Expected outcome if framework correct: the extracted material-culture and subsistence vocabulary overlaps the agriculture-and-technology native-vocabulary signatures identified in Section~\ref{sec:ch3} at high proportions (approximately the same as the 189-term curated kakawin set), and geographic references cluster in the volcanic-interior zones even in narratives set in the court register. Expected outcome if wrong: material-culture vocabulary is predominantly Sanskrit, or geographic references do not cluster in the volcanic interior.

\textbf{Prediction 6: VOC \textit{dagregister} NLP for colonial-era find reports.} An extension of the Channel 5 methodology to Dutch East India Company diary records (1602--1798) using contemporary computational-linguistics approaches to historical Dutch \citep{verberne2024}. Expected outcome if framework correct: colonial-era reports of ``oudheidkundig'' finds in Java cluster statistically in predicted-burial-depth zones at $p < 0.01$, replicating the Channel 5 pattern on a corpus independent of the 1854--1949 OV and Delpher sources used there. Expected outcome if wrong: no clustering or clustering away from predicted zones.

\textbf{Prediction 7: Phytolith and microbotanical extraction at Liangan or equivalent.} Liangan (Central Java) is a Hindu-Buddhist settlement buried under 9th-century Merapi ash and currently under active excavation. The prediction is extended to any similarly buried non-elite deposit. Expected outcome if framework correct: phytolith assemblages include domesticated rice and tuber signatures consistent with Hindu-Buddhist-era subsistence, and stratigraphic sequences below the 9th-century horizon contain older material-culture indicators. Expected outcome if wrong: absence of subsistence phytoliths or stratigraphic discontinuity suggesting isolated single-event occupation rather than continuous population.

\textbf{Prediction 8: DEM-based depression detection plus field verification.} A systematic digital elevation model analysis of Java's volcanic basins to identify anthropogenic-candidate depressions (deflations, artificial fill patterns, geometric regularities) below the noise threshold of natural volcanic-geomorphology variation. Expected outcome if framework correct: at least ten high-confidence candidate depressions detectable, with at least one returning archaeological material when verified via coring or test excavation. Expected outcome if wrong: no candidates above noise threshold, or no archaeological material in verified candidates.

\paragraph{Archaeometric predictions added in v0.3 (Channel 6).}

\textbf{Prediction 9:} Chemical fingerprinting of any future-excavated Javanese or Balinese pre-400~CE glass or bronze should match the Jia et al.\ (2024) and Calo (2014) compositional signatures at standard peer-review methodological thresholds. \textbf{Prediction 10:} The anticipated 10--20 unreported \textit{Jatim}-type finds in Korean, Japanese (Kofun period), or Chinese tombs currently awaiting archaeometric characterisation should, when analysed, return Java-consistent signatures under the full-signature diagnostic. \textbf{Prediction 11:} Systematic compositional study of pre-Tarumanagara West Javanese assemblages (Buni, Batujaya pre-temple layer) should detect trace-element signatures consistent with volcanic-interior provenance if the within-island trade inferred in Section~\ref{sec:ch2} is real.

\paragraph{Stop criteria at the framework level.}

The framework as a whole is considered falsified (as distinct from requiring revision of specific predictions) under any of the following outcomes:

\begin{enumerate}
\item Three or more of Predictions 1--8 fail independently, that is, yield outcomes consistent with the ``if framework wrong'' case rather than the ``if framework correct'' case.
\item The Philippines pre-400~CE site record is substantially retracted or re-dated to post-400~CE horizons in peer-reviewed re-analysis that supersedes the current consensus.
\item An alternative, successful peer-reviewed production-origin analysis relocates the Pejeng bronze drum tradition outside Bali/Java (removing the pre-inscription anchor of Channel 6).
\item Systematic East Javanese WGS results (Prediction 3) invert the polarity of the lineage-diversity and admixture-timing predictions.
\end{enumerate}

These criteria are specified in advance of their test. They are falsifiers rather than confirmation conditions: the framework's continued viability requires that none of them triggers, and the purpose of stating them explicitly is to establish what outcomes would cause the authors to abandon the present argument.


%% ================================================================
\section{Limitations}
\label{sec:limits}

The argument developed in this paper rests on evidence that is computational, interpretive, and comparative. It does not rest on new physical excavation. We catalogue its limitations explicitly.

\textit{Zero new field evidence.} Every claim in this paper derives from published data, pre-existing archaeological and linguistic corpora, or computational analyses of those resources. No core has been extracted; no site has been excavated; no artifact has been newly recovered. The framework is designed to generate specific predictions for future field work, and its evidential base will be strengthened or weakened by the outcomes of that work. As it stands, the argument is a synthesis, not a discovery.

\textit{Single research-group echo chamber.} Approximately 207 VOLCARCH experiments have been conducted by one human-computational dyad. The present paper has benefitted from cross-model critical review prior to submission (results documented in external review files), and has been revised in response to that review. Even so, peer engagement with the archaeological community has not yet been extensive, and the present synthesis is a stage in a research programme rather than its conclusion. The principal risk is that systematic blind spots in the computational approach remain unrecognised; the principal mitigation is explicit acknowledgement of this risk and invitation to the archaeological community to test it.

\textit{Cascade underdetermination.} As noted in Section~\ref{sec:layers}, multiple parameterisations of a six-filter multiplicative visibility model can bracket the observed archaeological record. The cascade is diagnostic (identifying where the greatest leverage lies) rather than predictive (allowing any specific observation to be generated from the filter set). The pattern of convergence across the six channels does not depend on the cascade; it is the cascade that depends on the pattern.

\textit{Dataset concentration in DHARMA.} Approximately twenty-five VOLCARCH experiments depend on the 268 Javanese inscriptions catalogued in DHARMA. The Channel 5 colonial-archive corpus, Channel 6 archaeometric literature, Channel 4 paleogenomic reinterpretation, and Channel 3 Austronesian Comparative Dictionary material mitigate this monoculture but do not eliminate it. A reader sceptical of DHARMA-based inferences would see the present argument as more than one-quarter supported by a single corpus type.

\textit{Limits of comparative independence.} The Philippines comparison (Section~\ref{sec:ch2}) shares Austronesian cultural heritage with Java; the comparison controls for tropical ecology and organic architecture but not for all possible confounders. The West Java within-island control (Buni, Batujaya) is more fully matched but samples only coastal contexts; its extension inland through the Bekasi--Jakarta transect would be required for full control (Prediction 4).

\textit{Deep-interior claim is about demographic magnitude, not archaeometric provenance.} One reviewer of an earlier draft correctly noted that the Channel 6 archaeometric evidence (Pejeng drums, \textit{Jatim} beads) does not prove production in the deep volcanic interior specifically. The archaeometric channel establishes that production occurred \textit{somewhere} in Java and Bali and that the workshops are archaeologically absent from the entire Java--Bali region. The deep-interior claim of the present paper rests on Channels 1 and 2 (sedimentation plus Philippines comparison), which specifically concern the volcanic interior. Channel 6's contribution to the deep-interior claim is indirect: it establishes that Javanese production systems have left archaeological traces outside Java and Bali but not within them, a pattern consistent with the filter set developed here but not itself proof of any particular production location.

\textit{The Gunung Padang precedent.} The retracted paper by \citet{natawidjaja2023} claimed very old radiocarbon dates for putative artificial structures at Gunung Padang in West Java. The retraction was appropriate: radiocarbon dating of charcoal at depth is insufficient to establish archaeological presence without diagnostic cultural material. We note this case as an object lesson in the present framework. The predictions of Section~\ref{sec:predictions} require \textit{both} stratigraphic dating \textit{and} diagnostic artifacts; the framework commits to the same standard the Gunung Padang retraction enforced. Any future claim of pre-Hindu deposits in Java's volcanic interior that rests on dating alone, without artifact recovery, will not support this framework and should not be treated as doing so.

\textit{Selective-survival framing risks unfalsifiability without discipline.} As acknowledged at the close of Section~\ref{sec:selective_survival}, the selective-survival framing could, without discipline, absorb any outcome. The framework is disciplined by the specific predictions and thresholds of Section~\ref{sec:predictions}. The framework is not the hypothesis; the predictions and their pre-registered thresholds are.


%% ================================================================
\section{Conclusions}
\label{sec:conclusion}

The archaeological record of Java before the fourth century~CE is effectively empty. We have argued, through the convergence of six independent evidence channels, that the emptiness is a consequence of compound taphonomic processes operating on a substantial substrate civilisation rather than evidence that the civilisation itself was absent or demographically trivial. The pre-400~CE population of Java was, on three independent demographic estimation methods, on the order of one to two million inhabitants, with a conservative lower bound of approximately 600{,}000 and a very conservative lower bound (under doubled post-400~CE growth rates) of approximately 150{,}000--250{,}000. Even under the most conservative baseline the expected-versus-observed settlement gap exceeds two orders of magnitude.

The argument does not rest on any single line. Sedimentation rate calibration places pre-Hindu deposits at 4--10~m depth in volcanic Java, beyond the reach of standard surface survey and conventional GPR. Deep comparison with the Philippine archaeological record documents 275--340 pre-400~CE sites in a neighbouring Austronesian region with lower volcanic density and more abundant karst, against zero in Java's volcanic interior. Linguistic substrate reconstruction identifies Proto-Austronesian writing vocabulary, hundreds of substrate lexical items in modern Austronesian languages, and a domain gradient in Old Javanese literary sources that is compatible with a pre-existing indigenous lexical substrate for agriculture, landscape, and cosmology. Paleogenomic reinterpretation documents deep lineage-specific variation in modern West Javanese WGS and an archipelago-wide admixture pattern that cannot be accounted for by recent migration alone; the absence of ancient DNA from volcanic Java is acknowledged to be a downstream consequence of the absence of skeletal remains rather than independent evidence of a separate preservation filter. Text mining of Dutch colonial archives surfaces volcano-distance gradients and depth distributions consistent with the sedimentation model. Archaeometric evidence in two chronological layers---1st--3rd century~CE Pejeng bronze drums locally produced in Bali and Java, and 4th--6th century~CE \textit{Jatim} glass beads in elite tombs from Northern Wei China to Silla Korea to Roman Berenike---documents an invisibility regime that operates on organised production infrastructure across six centuries, from the pre-inscription period through the early state-formation period.

Wayang theatre preserves, in a living tradition, approximately 20--30\% of its narrative repertoire in a sempalan layer with no Indian source, including the Punakawan clown-servant complex whose cosmological role and etymology are Austronesian and whose pre-Hindu dating is confirmed by the 1358~CE Candi Tigamangi relief. The residue of pre-Hindu cultural material in the living tradition is not proof of a pre-Hindu civilisation; it is evidence that pre-Hindu cultural content was sufficient in volume and structure to survive, measurably, more than a millennium of cultural layering.

The framework organising these findings is a six-filter erasure model (volcanic burial, coastal submersion, historiographic bias, cosmological overwrite, genre taphonomy, historiographic periodicity), with karst availability as a seventh modulating factor. The filters combine multiplicatively; the severity of the volcanic-Java case derives from their joint rather than any single operation. The framework is not predictive in the strong sense: multiple parameterisations can bracket the observed record. It is diagnostic: it identifies where leverage lies and where remedial field research would have the greatest payoff.

Eight pre-registered predictions with specified falsification thresholds---plus three archaeometric sub-predictions---commit the framework to outcomes that can confirm, constrain, or falsify its central claim. The framework is considered falsified if three or more predictions fail, if the Philippines comparison is substantially retracted, if the Pejeng bronze tradition is relocated outside Java and Bali, or if East Javanese whole-genome sequencing inverts the framework-predicted bottleneck signature.

The methodological implication of this argument is broader than any specific finding about Java's pre-Hindu period. Where taphonomy is severe and materially selective, single-channel archaeology becomes insufficient. Recovering the record of such settings requires multi-channel convergence analysis, combined with the specific field methods---deep coring, archaeometric provenance analysis, paleogenomic contextualisation, historical-text NLP, living-tradition analysis---that can access the register in which the record has survived. The framework developed here is a general-purpose approach to invisibility problems in archaeology, of which volcanic tropical Island Southeast Asia is one instance among several.

Indonesia's archaeological narrative has long been structured around the appearance of Hindu-Buddhist kingdoms and their monumental record. The framework developed here does not contest the reality of that appearance; it contests the interpretation of what the appearance replaced. Pre-Hindu Nusantara was not absent. It was archaeologically invisible---and it still is, outside the channels assembled here. Making it visible, materially, is the research programme the framework points toward.

\bibliographystyle{apalike}
\bibliography{references}

\end{document}
